\scanpage{34}
Soit $G$ un \inserted{\unsure{schéma en}} groupe affine sur $A$, $A$-\unsure{simple}
\[
\begin{array}{ll}
\operatorname{Aff}_A(G)&\text{son algèbre affine}\\[1ex]
U_A(G)&\text{son hyperalgèbre}
\end{array}
\]

Donc
\[
\left\{\begin{array}{l}
\operatorname{Aff}_A(G)\text{ est une }A\text{-algèbre commutative augmentée,}\\
\qquad\text{munie d'un hom diagonal multiplicatif,}\\
\qquad\text{associatif, compatible avec les augmentations}\\[1ex]
U_A(G)\text{ est une }A\text{-algèbre en général non comm.,}\\
\qquad\text{munie d'un hom diagonal multiplicatif,}\\
\qquad\text{associatif, comm. compatible avec les}\\
\qquad\text{augmentations}
\end{array}\right.
\]

Il y a un produit scalaire
\[
\langle f,X\rangle=\langle\deleted{\unsure{f}}\inserted{X},\deleted{\unsure{X}}\inserted{f}\rangle
\]
entre $\operatorname{Aff}_A(G)$ et $U_A(G)$, ce dernier s'identifiant aux formes $A$-linéaires sur $\operatorname{Aff}_A(G)$ qui s'annulent sur une puissance de l'idéal d'augmentation.

$U_A(G)$ opère à g. et à droite sur $\operatorname{Aff}_A(G)$ [qui est un bimodule sur $U_A(G)$]
\[
\left\{\begin{aligned}
\langle Y,X*f\rangle&=\langle Y*X,f\rangle=\langle f,Y*X\rangle=\langle f*Y,X\rangle\\
\langle\varepsilon,X*f\rangle&=\langle\varepsilon,f*X\rangle=\langle X,f\rangle
\end{aligned}\right.
\]
\margin{33}
